
\section*{Methods}

\subsection*{Data sources}

\textit{Biodiversity data:} All biodiversity data -- species presence or absence, and abundance -- were obtained from the Projecting Responses of Ecological Diversity in Changing Terrestrial Systems (PREDICTS) project \cite{hudson2017}, a meta-database compiled from independent source studies and inventories. We combined the two publicly available releases of the database (from 2016 and 2022) \cite{hudson2023,contu2022} into one dataset. Before filtering (see further down), the combined dataset contained data from 817 studies comprising 35,736 sampling sites in 101 countries, with 4,318,808 unique records across 53,925 species, collected between 1984 and 2018. The PREDICTS data mainly covers animals and plants, with the most common groups being insects and other arthropods, vascular plants, and vertebrate animals. Despite efforts to balance the data taxonomically, vertebrates are over-represented relative to their true diversity, while fungi and non-arthropod invertebrates are underrepresented. Geographically, the dataset has gaps in line with most biodiversity data sources; high-income regions like North America and Europe are well-sampled relative to their underlying biodiversity, some tropical areas in Latin America also have high coverage, while there are significant gaps in Africa and parts of Asia (see \autoref{fig:predicts_data} for an overview of PREDICTS data coverage). In 255 of the source studies, spatially adjacent sampling sites had been grouped into spatial blocks, which was utilized for model training.

\textit{Human pressure data:} The predominant land-use type and land-use intensity at each sampling site has been categorized by the PREDICTS team based on information in the source studies \cite{hudson2017}. Land-use categories include primary vegetation, secondary vegetation (split into young, intermediate, mature, or indeterminate age), plantation forest, cropland, pasture and urban, while use intensity has been classified as minimal, light or intense. Data on human population density, expressed as the number of people per 1~km$^2$, came from the Gridded Population of the World, v4.11 (GPWv4.11) dataset \cite{gpw2017}, based on the 2010 round of Population and Housing Censuses (conducted 2005--2014) and adjusted to match the United Nations World Population Prospects country totals. The population data represents extrapolated numbers for the years 2000, 2005, 2010, 2015, and 2020.  Data on road networks were taken from the Global Roads Open Access Data Set, v1 (gROADS) \cite{groads2013}, a global layer of joined country road networks adjusted for topology. Data were collected between 1980 and 2010, with limited information on original road construction dates.

\textit{Bioclimatic and topographic data:} Bioclimatic variables, such as annual averages, seasonality and extreme values of temperature and precipitation, were based on the WorldClim v2.1 dataset\cite{fick2017}, frequently used in species distribution models and other biodiversity studies. The data represent average values over the period 1970--2000. Data on elevation, slope and terrain roughness came from the EarthEnv repository \cite{amatulli2018}, constructed from global digital elevation models derived from satellite imagery and LiDAR measurements.

The spatial resolution of the PREDICTS land-use data depends on the specific sampling extent of a given site in each source study \cite{hudson2017} (which is only available for 13.1\% of sampling sites). The resolution of the population density, bioclimatic and topographic variables is 30 arc-seconds (approximately 1~km$^2$ at the equator) \cite{gpw2017,fick2017, amatulli2018}. The spatial accuracy of the road network data varies by country\cite{groads2013}. Manual classification of land-use type and intensity could have introduced inconsistencies within and between studies \cite{hudson2017}. Since predominant land-use is a categorical attribute, there might be underlying differences in habitat conditions within the same class that are notn observed. The remote sensing and census data layers are all modeled to some extent, potentially introducing additional uncertainty and errors.
\vspace{\baselineskip}
